\section{Node-Accurate Area Calculation}
\label{sec:primfac}
In node-accurate area calculation based on prime factor encoding, each sector of the stator and rotor is assigned a unique prime number, where \(\gls{Primenumber}_{k}\) denotes the \(k\)-th prime number. The results of this assignment are depicted in Fig. \ref{fig:Primfaktor_Stator} for the stator and with consecutive encoding in Fig. \ref{fig:Primfaktor_Rotor} for the rotor.
%
\begin{figure}[ht]%
    \centering%
    \begin{subfigure}{.49\linewidth}%
        \centering%
        \begin{tikzpicture}[font=\footnotesize]
    \pgfmathsetmacro{\slots}{3}
    \pgfmathsetmacro{\innerRadius}{1.2}
    \pgfmathsetmacro{\outerRadius}{2.2}
    \pgfmathsetmacro{\meanRadius}{0.5*(\innerRadius + \outerRadius)}
    \pgfmathsetmacro{\segmentangle}{360/\slots}
    \pgfmathsetmacro{\slotWidth}{30}
    \pgfmathsetmacro{\shoeWidth}{30}
    
    \pgfmathsetmacro{\innerSlotwidth}{\slotWidth/\innerRadius}
    \pgfmathsetmacro{\outerSlotwidth}{\slotWidth/\outerRadius}
    \pgfmathsetmacro{\innerShoewidth}{\shoeWidth/\innerRadius}
    \pgfmathsetmacro{\outerShoewidth}{\shoeWidth/\outerRadius}
    \pgfmathsetmacro{\meanShoewidth}{0.5*(\innerShoewidth + \outerShoewidth)}
    
    \clip (0,0) circle(\outerRadius);

    \pgfmathsetmacro{\iters}{\slots-1}
    \foreach \k in {0,...,\iters}
    {
        \pgfmathsetmacro{\innerCoreCW}{\k*\segmentangle - (\segmentangle - \innerSlotwidth)/2 + \innerShoewidth}
        \pgfmathsetmacro{\innerCoreCCW}{\k*\segmentangle + (\segmentangle - \innerSlotwidth)/2 -\innerShoewidth}
        \pgfmathsetmacro{\outerCoreCW}{\k*\segmentangle - (\segmentangle - \outerSlotwidth)/2 + \outerShoewidth}
        \pgfmathsetmacro{\outerCoreCCW}{\k*\segmentangle + (\segmentangle - \outerSlotwidth)/2 - \outerShoewidth}
    
        \filldraw[fill=colIron] 
        (\innerCoreCW:\innerRadius) arc[start angle=\innerCoreCW, end angle=\innerCoreCCW, radius=\innerRadius] --
        (\outerCoreCCW:\outerRadius) arc[start angle=\outerCoreCCW, end angle=\outerCoreCW, radius=\outerRadius] -- 
        cycle;

        \filldraw[fill=colIronLight] 
        (\innerCoreCW:\innerRadius) arc[start angle=\innerCoreCW, end angle={\innerCoreCW - \innerShoewidth}, radius=\innerRadius] --
        ({\outerCoreCW - \outerShoewidth}:\outerRadius) arc[start angle={\outerCoreCW - \outerShoewidth}, end angle=\outerCoreCW, radius=\outerRadius] -- 
        cycle;

        \filldraw[fill=colIronLight] 
        (\innerCoreCCW:\innerRadius) arc[start angle=\innerCoreCCW, end angle={\innerCoreCCW + \innerShoewidth}, radius=\innerRadius] --
        ({\outerCoreCCW + \outerShoewidth}:\outerRadius) arc[start angle={\outerCoreCCW + \outerShoewidth}, end angle=\outerCoreCCW, radius=\outerRadius] -- 
        cycle;

        \draw (0,0) circle[radius=\innerRadius];
        \draw (0,0) circle[radius=\outerRadius];

        \pgfmathsetmacro{\idxGap}{int(1 + \k*4)}
        \pgfmathsetmacro{\idxShoeCW}{int(2 + \k*4)}
        \pgfmathsetmacro{\idxCore}{int(3 + \k*4)}
        \pgfmathsetmacro{\idxShoeCCW}{int(4 + \k*4)}        
        
        \path ({(\k-0.5)*\segmentangle}:\innerRadius) -- ({(\k-0.5)*\segmentangle}:\outerRadius) node[pos=0.5, sloped] {$\gls{Primenumber}_{\idxGap}$};
        \path ({\k*\segmentangle}:\innerRadius) -- ({\k*\segmentangle}:\outerRadius) node[pos=0.5, sloped] {$\gls{Primenumber}_{\idxCore}$}; 
        
        \path ({\innerCoreCW - 0.5*\innerShoewidth}:\innerRadius) -- ({\outerCoreCW - 0.5*\outerShoewidth}:\outerRadius) node[pos=0.5, sloped] {$\gls{Primenumber}_{\idxShoeCW}$};     
        
        \path ({\innerCoreCCW + 0.5*\innerShoewidth}:\innerRadius) -- ({\outerCoreCCW + 0.5*\outerShoewidth}:\outerRadius) node[pos=0.5, sloped] {$\gls{Primenumber}_{\idxShoeCCW}$};
    }    
    
\end{tikzpicture}%
        \caption{Stator}%
        \label{fig:Primfaktor_Stator}%
    \end{subfigure}%
    \hfill%
    \begin{subfigure}{.49\linewidth}%   
        \centering%
        \begin{tikzpicture}[font=\footnotesize]
    \pgfmathsetmacro{\poles}{2}
    \pgfmathsetmacro{\innerRadius}{1.2}
    \pgfmathsetmacro{\outerRadius}{2.2}
    \pgfmathsetmacro{\meanRadius}{0.5*(\innerRadius + \outerRadius)}

    \pgfmathsetmacro{\mcf}{0.8}
    \pgfmathsetmacro{\poleWidth}{360/\poles}
    \pgfmathsetmacro{\magnetWidth}{\poleWidth*\mcf}
    \pgfmathsetmacro{\gapWidth}{\poleWidth*(1-\mcf)}  

    \clip (0,0) circle(\outerRadius);
    
    \filldraw[fill=colMagnetNorth] 
    ({\magnetWidth/2}:\innerRadius) arc[start angle={\magnetWidth/2}, end angle={-\magnetWidth/2}, radius=\innerRadius] --
    ({-\magnetWidth/2}:\outerRadius) arc[start angle={-\magnetWidth/2}, end angle={\magnetWidth/2}, radius=\outerRadius] -- 
    cycle;

    \filldraw[fill=colMagnetSouth] 
    ({\magnetWidth/2 + \poleWidth}:\innerRadius) arc[start angle={\magnetWidth/2 + \poleWidth}, end angle={-\magnetWidth/2 + \poleWidth}, radius=\innerRadius] --
    ({-\magnetWidth/2 + \poleWidth}:\outerRadius) arc[start angle={-\magnetWidth/2 + \poleWidth}, end angle={\magnetWidth/2 + \poleWidth}, radius=\outerRadius] -- 
    cycle;

    \draw (0,0) circle[radius=\innerRadius];
    \draw (0,0) circle[radius=\outerRadius];     
    
    \path (0,0) -- (0:\meanRadius) node[pos=1, sloped] {$\gls{Primenumber}_{13}$};
    \path (0,0) -- (90:\meanRadius) node[pos=1, sloped] {$\gls{Primenumber}_{14}$};
    \path (0,0) -- (180:\meanRadius) node[pos=1, sloped] {$\gls{Primenumber}_{15}$};        
    \path (0,0) -- (270:\meanRadius) node[pos=1, sloped] {$\gls{Primenumber}_{16}$};    
\end{tikzpicture}%
        \caption{Rotor}%
        \label{fig:Primfaktor_Rotor}%
    \end{subfigure}%
    \caption{Exemplary prime factor encoding for a machine with \(\gls{Ntooth} = 3\) and \(\gls{polepair} = 1\)}%
    \label{fig:Primfactor_Codierung}%
\end{figure}%
%
Next, vectors are defined for both the stator and rotor, with the number of elements in each vector proportional to the geometric width of each sector; wider sectors correspond to a higher number of elements associated with the \(k\)-th prime number. To illustrate this for the rotor: 
%
$
    \gls{vec_prim_sec}_{\gls{idx_rotor}} =
    \begin{bsmallmatrix}
        \gls{Primenumber}_{13} &
        \gls{Primenumber}_{13} &
        \gls{Primenumber}_{13} &
        \gls{Primenumber}_{13} &
        \gls{Primenumber}_{13} &
        \gls{Primenumber}_{14} &
        \gls{Primenumber}_{15} &
        \gls{Primenumber}_{15} &
        \gls{Primenumber}_{15} &
        \gls{Primenumber}_{15} &
        \gls{Primenumber}_{15} &
        \gls{Primenumber}_{16}
    \end{bsmallmatrix}^{\gls{trans}}.
$
%
A vector \(\gls{vec_prim_sec}_{\gls{idx_stator}}\) of equal length is defined for the stator, constructed following the same principle. Element-wise multiplication \( \gls{vec_prim_sec}_{\gls{schurr}} = \gls{vec_prim_sec}_{\gls{idx_rotor}} \, \gls{schurr} \, \gls{vec_prim_sec}_{\gls{idx_stator}}\) of these vectors yields a result vector representing the intersection areas and can be interpreted using a result matrix \(\gls{matr_prim_res}\). To construct the result matrix, the prime factors for the stator
$
\gls{vec_prim}_{\gls{idx_stator}} 
=
\begin{bsmallmatrix}
    \gls{Primenumber}_{1} & \cdots & \gls{Primenumber}_{12}
\end{bsmallmatrix}^{\gls{trans}}
$ 
and rotor 
$
\gls{vec_prim}_{\gls{idx_rotor}} 
=
\begin{bsmallmatrix}
    \gls{Primenumber}_{13} & \cdots & \gls{Primenumber}_{16}
\end{bsmallmatrix}^{\gls{trans}}
$
are encoded in vectors and then multiplied dyadically:
%
\begin{equation}
        \gls{matr_prim_res} =
        \gls{vec_prim}_{\gls{idx_rotor}} \cdot \gls{vec_prim}_{\gls{idx_stator}}^{\gls{trans}} =
        \begin{bsmallmatrix}
            82 & 123 & 205 & 287 & 451 & 533 & 697 & 779 & 943 & 1189 & 1271 & 1517 \\
            86 & 129 & 215 & 301 & 473 & 559 & 731 & 817 & 989 & 1247 & 1333 & 1591 \\
            94 & 141 & 235 & 329 & 517 & 611 & 799 & 893 & 1081 & 1363 & 1457 & 1739 \\
            106 & 159 & 265 & 371 & 583 & 689 & 901 & 1007 & 1219 & 1537 & 1643 & 1961
        \end{bsmallmatrix}
        \label{equ:PrimResMatrix}.
\end{equation}
%
In $\gls{matr_prim_res}$ the row corresponds to the rotor sector and the column to the stator sector. The count of \(\gls{Primenumber}_{k}\) in $\gls{vec_prim_sec}_{\gls{schurr}}$ is related to the area of the sectors. To account for rotor position dependence, \(\gls{vec_prim_sec}_{\gls{idx_rotor}}\) is cyclically shifted. The calculation of \(\gls{vec_prim_sec}_{\gls{schurr}}\) is then repeated with the shifted vector and interpreted using \(\gls{matr_prim_res}\) to derive the rotor position-dependent areas. For an exemplary operating point, this is depicted in Fig. \ref{fig:Primfaktor_Result}.
%
\begin{figure}[ht]
    \centering
    \begin{tikzpicture}
    \def\filename{figures/MEC/Primfactor_Area.csv}
    \begin{groupplot}[
        group style={group size=3 by 1, horizontal sep=1.5cm},
        width=\linewidth/3,
        height=3.5cm,
        xlabel={$\gls{idx_dim_azimutal}/\si{\degree}$},
        ylabel={},
        grid=both,
        grid style={dashed},
        xmin = 0, 
        xmax = 360,
        xtick={0,120,...,360},
    ]

    \nextgroupplot[title = {$\gls{vec_prim_sec}_{\gls{idx_stator}}$}]
    \addplot[const plot, colOne, line width = 1] table [x expr = {\thisrow{theta}}, y = Stator] {\filename};   

    \nextgroupplot[title = {$\gls{vec_prim_sec}_{\gls{idx_rotor}}$}]
    \addplot[const plot, colOne, line width = 1] table [x expr = {\thisrow{theta}}, y = Rotor] {\filename};

    \nextgroupplot[title = {$\gls{vec_prim_sec}_{\gls{schurr}}$}]
    \addplot[const plot, colOne, line width = 1] table [x expr = {\thisrow{theta}}, y = Intersect] {\filename}; 
    \addplot [draw = none, fill=colTwo, fill opacity=0.3, no marks] coordinates {(161,0) (160,901) (195,901) (195,0) (160,0)};

    \end{groupplot}
\end{tikzpicture}

    \caption{Exemplary result of area calculation using prime factor encoding for a \(\gls{Ntooth} = 3\), \(\gls{polepair}=1\) machine}
    \label{fig:Primfaktor_Result}
\end{figure}
%
The marked area corresponds to the value \(\gls{vec_prim_sec}_{\gls{schurr}} = 901\). From (\ref{equ:PrimResMatrix}), it follows that this element is located at position \(\gls{matr_prim_res}(4,7)\), indicating a cross-sectional area between rotor sector 4 and stator sector 7. Furthermore, the sector width is \(\Delta \gls{idx_dim_azimutal}\), and its area can be determined as \(\gls{area} = \frac{\pi}{\SI{180}{\degree}} \cdot \Delta \gls{idx_dim_azimutal} \cdot \gls{meanRadius} \cdot (\gls{outerRadius} - \gls{innerRadius})\).
